\documentclass[11pt]{article}
\usepackage[margin=1in]{geometry}
\usepackage{amsmath,amssymb,amsthm,microtype}
\usepackage[colorlinks=true,linkcolor=blue,urlcolor=blue]{hyperref}
\newtheorem{lemma}{Lemma}
\newtheorem{theorem}{Theorem}
\newtheorem{corollary}{Corollary}
\newcommand{\dist}{\operatorname{dist}}
\newcommand{\np}[2]{\lVert#1\rVert_{#2}}
\title{A Summable Budget for Adaptive Primal Weights\\\large Companion to the Mapped-Point Restart Proof}
\author{}
\date{November 11, 2026}
\begin{document}
\maketitle

This standalone review note concerns the experimental weight clip, not the unchanged production controller.
It uses the restricted KL-UOT KKT mapping $\mathcal M$ and solution set $\mathcal S$ from \texttt{mapped\_restart\_convergence.tex}.
That note establishes existence of a KKT point, firm nonexpansiveness of each fixed-metric PDHG map, the Halpern epoch estimates, and the local KKT error bound.
All proximal evaluations and inequalities below are in exact arithmetic.

\section{Weight proposal and metric comparison}

Let $\overline{\mathbf D}_x,\overline{\mathbf D}_y$ be fixed positive diagonal base steps satisfying $\zeta:=\|\overline{\mathbf D}_y^{1/2}\mathbf A\overline{\mathbf D}_x^{1/2}\|_2<1$.
For a log weight $s=\log w$, write
\begin{equation}
\mathbf P(s)=\begin{pmatrix}e^{-s}\overline{\mathbf D}_x^{-1}&-\mathbf A^\top\\-\mathbf A&e^s\overline{\mathbf D}_y^{-1}\end{pmatrix},\qquad T_s=(\mathbf P(s)+\mathcal M)^{-1}\mathbf P(s).
\end{equation}
Within epoch $j$, keep $s_j$ fixed and perform the reflected Halpern iteration from anchor $\mathbf z_j^{\mathrm a}$.
For any finite epoch length $L_j\geq1$, set $\mathbf z_{j+1}^{\mathrm a}=T_{s_j}\mathbf z^{j,L_j-1}$.
Then obtain a proposed next weight $\widetilde w_{j+1}>0$ from the existing controller, including its fallback branch, and use
\begin{equation}
s_{j+1}=\Pi_{[\log w_{\min},\log w_{\max}]}\!\left(s_j+\operatorname{clip}(\log\widetilde w_{j+1}-s_j,-b_j,b_j)\right),\qquad b_j=\frac{B}{(j+1)^p},\quad B>0,\quad p>1.
\label{eq:clip}
\end{equation}
Here $0<w_{\min}\leq w_0\leq w_{\max}<\infty$.
The final interval projection does not increase the displacement from $s_j$.
The index counts every restart within one restricted call and resets only at a new restricted problem.
There is no restriction on how the proposal is generated.
If its original update lies within the budget, an implementation can return the original floating-point weight directly, avoiding an unnecessary log--exp round trip.

\begin{lemma}[Controlled change of metric]
For any real $s,t$,
\begin{equation}
e^{-c|t-s|}\mathbf P(s)\preceq\mathbf P(t)\preceq e^{c|t-s|}\mathbf P(s),\qquad c=(1-\zeta)^{-1}.
\label{eq:comparison}
\end{equation}
In particular, $\mathbf P_j=\mathbf P(s_j)$ converges to a positive definite matrix, is uniformly bounded and uniformly positive definite, and satisfies
\begin{equation}
\mathbf P_{j+1}\preceq r_j\mathbf P_j,\qquad r_j=e^{cb_j},\qquad \prod_{j=0}^{\infty}r_j<\infty.
\label{eq:product}
\end{equation}
\end{lemma}

\begin{proof}
Let $\mathbf D(s)$ be the block-diagonal part of $\mathbf P(s)$.
The scaled off-diagonal block has norm $\zeta$, independently of $s$, so $\mathbf P(s)\succeq(1-\zeta)\mathbf D(s)$.
For any nonzero $\mathbf h$,
\begin{equation}
\left|\frac{d}{ds}\log(\mathbf h^\top\mathbf P(s)\mathbf h)\right|\leq\frac{\mathbf h^\top\mathbf D(s)\mathbf h}{\mathbf h^\top\mathbf P(s)\mathbf h}\leq c.
\end{equation}
Integration gives~\eqref{eq:comparison}.
Equation~\eqref{eq:clip} gives $|s_{j+1}-s_j|\leq b_j$ and
\begin{equation}
\sum_{j=0}^{\infty}|s_{j+1}-s_j|\leq\sum_{j=0}^{\infty}b_j\leq B\left(1+\frac{1}{p-1}\right)<\infty.
\end{equation}
Thus $s_j$ converges, and the compact weight interval gives the uniform spectral bounds and positive definiteness of the limit.
The product bound follows from $\sum_jb_j<\infty$.
\end{proof}

\section{Global convergence with the clipped controller}

\begin{theorem}[Global convergence]
For arbitrary controller proposals and arbitrary mapped-point restart times, the clipped-weight iteration converges to a single restricted KKT point.
This includes anchors, chronological Halpern iterates, mapped points, and the primal projection multipliers.
\end{theorem}

\begin{proof}
First suppose there are infinitely many finite epochs.
Let $\mathbf t_j=\mathbf z^{j,L_j-1}$, $e_j=\np{\mathbf t_j-T_{s_j}\mathbf t_j}{\mathbf P_j}$, and fix $\mathbf z^\star\in\mathcal S$.
The firm nonexpansiveness and epoch-ball estimates in the companion note give
\begin{equation}
\np{\mathbf z_{j+1}^{\mathrm a}-\mathbf z^\star}{\mathbf P_j}^{2}+e_j^2\leq\np{\mathbf z_j^{\mathrm a}-\mathbf z^\star}{\mathbf P_j}^{2}.
\end{equation}
Set $E_j=\np{\mathbf z_j^{\mathrm a}-\mathbf z^\star}{\mathbf P_j}^{2}$ and $\Gamma_j=\prod_{i=0}^{j-1}r_i$.
The metric comparison yields
\begin{equation}
E_{j+1}\leq r_j(E_j-e_j^2),\qquad \frac{E_{j+1}}{\Gamma_{j+1}}\leq\frac{E_j}{\Gamma_j}-\frac{e_j^2}{\Gamma_j}.
\label{eq:scaled}
\end{equation}
Since $1\leq\Gamma_j\leq\Gamma_\infty<\infty$, the anchors are bounded, $E_j$ has a limit, and $\sum_je_j^2<\infty$.
The resolvent identity gives
\begin{equation}
\mathbf P_j(\mathbf t_j-\mathbf z_{j+1}^{\mathrm a})\in\mathcal M(\mathbf z_{j+1}^{\mathrm a}).
\end{equation}
Its left-hand side tends to zero by uniform spectral bounds and $e_j\to0$.
Every cluster point of the anchors therefore belongs to $\mathcal S$ by closedness of the maximal monotone graph.
Choose one cluster point $\mathbf z^\infty\in\mathcal S$.
The preceding argument shows that $E_j$ relative to this point has a limit, and its convergent subsequence makes that limit zero.
Uniform positive definiteness implies $\mathbf z_j^{\mathrm a}\to\mathbf z^\infty$.
Within each epoch, both the Halpern and mapped iterates have $\mathbf P_j$-distance to $\mathbf z^\infty$ at most that of the anchor.
Uniform spectral bounds prove their chronological convergence as well.
The projection multiplier is
\begin{equation}
\boldsymbol\lambda^{j,k}=\mathbf c+\mathbf A^\top\mathbf y^{j,k}-\mathbf D_{x,j}^{-1}(\mathbf x^{j,k}-\widehat{\mathbf x}^{j,k}).
\end{equation}
It converges to $\mathbf c+\mathbf A^\top\mathbf y^\infty$ because $\mathbf D_{x,j}^{-1}$ is uniformly bounded.
If there are finitely many restarts, the weight is fixed in the final infinite epoch and the usual Halpern theorem supplies the same conclusions.
\end{proof}

\section{Tail linear rates}

\begin{theorem}[Linear rates under the existing local KL-UOT error bound]
The anchors converge R-linearly when indexed by epoch, provided infinitely many restarts occur.
If the sufficient-residual-reduction restart test is checked with bounded gaps and its baseline is refreshed under the new metric, convergence is also R-linear in primary PDHG evaluations after a finite initial phase.
The necessary-reduction and artificial-length triggers may remain active.
\end{theorem}

\begin{proof}
The local KKT error bound in the companion note depends on $\mathcal M$, not the weight.
Since the metrics have uniform spectral bounds, its transfer to the family of fixed-point maps gives a common local constant $\alpha_0>0$ such that
\begin{equation}
\alpha_0\dist_{\mathbf P_j}(\mathbf z,\mathcal S)\leq\np{\mathbf z-T_{s_j}\mathbf z}{\mathbf P_j}
\end{equation}
for all late epochs and all $\mathbf z$ sufficiently near $\mathbf z^\infty$.
For completeness, if the KKT error-bound constant is $C_{\mathrm K}$ and $\lambda_{\max}(\mathbf P_j)\leq\overline p$, one may take $\alpha_0=(1+C_{\mathrm K}\overline p)^{-1}$ after shrinking the neighborhood.
The mapped images stay in the required KKT neighborhood uniformly because every $T_{s_j}$ fixes $\mathbf z^\infty$ and is nonexpansive in a uniformly equivalent norm.
Let $d_j=\dist_{\mathbf P_j}(\mathbf z_j^{\mathrm a},\mathcal S)$ and $q_0=(1+\alpha_0^2)^{-1/2}$.
The fixed-metric epoch-contraction argument, followed by~\eqref{eq:comparison}, gives
\begin{equation}
d_{j+1}\leq e^{cb_j/2}q_0d_j.
\end{equation}
Since $b_j\to0$, choose $q_1\in(q_0,1)$ and $J$ such that $e^{cb_j/2}q_0\leq q_1$ for $j\geq J$.
Hence $d_j$ decays geometrically.
If $\lambda_{\min}(\mathbf P_j)\geq\underline p>0$, the Euclidean anchor increment is at most $2d_j/\sqrt{\underline p}$, so summing increments proves R-linear convergence to the actual limit, also uniformly within each epoch.
For sufficient-reduction factor $\theta\in(0,1)$ and check gaps at most $h$, the Halpern residual bound gives an eventual epoch-length bound $H=\lceil1/(\alpha_0\theta)\rceil+h$, exactly as in the companion note.
Consequently the epoch rate implies a bound $Cq_1^{\lfloor N/H\rfloor}$ in primary evaluations.
\end{proof}

\paragraph{What empirical trajectories can establish.}
A finite trace can identify a budget that leaves observed updates unchanged; it cannot prove summability of the unmodified infinite trajectory.
The intervention~\eqref{eq:clip} enforces the property on all possible trajectories, including fallback updates, so its guarantee is independent of the calibration sample.
Selecting a large $B$ may preserve current performance but makes the worst-case constants conservative.
The proof concerns convergence rather than a practically sharp complexity estimate.
There is no forced freezing time and no dependence on a safeguard rejection to activate the budget.

\end{document}
